\documentclass[11pt,a4paper]{article}
\usepackage[margin=2.5cm]{geometry}
\usepackage{xcolor}
\usepackage{titlesec}
\usepackage{enumitem}
\usepackage{fancyhdr}
\usepackage{hyperref}
\usepackage{parskip}

% PJD Brand Colours
\definecolor{pjdnavy}{HTML}{0A192F}
\definecolor{pjdgold}{HTML}{FFD700}

% Section formatting
\titleformat{\section}{\large\bfseries\color{pjdnavy}}{}{0em}{}[\color{pjdgold}\titlerule]
\titleformat{\subsection}{\normalsize\bfseries\color{pjdnavy}}{}{0em}{}

% Header/Footer
\pagestyle{fancy}
\fancyhf{}
\fancyhead[L]{\small\color{pjdnavy}\textbf{Paul Jeggels Designs}}
\fancyhead[R]{\small\color{pjdnavy}Surfboard Selection \& Customisation FAQ}
\fancyfoot[C]{\small\color{pjdnavy}15 Dageraad Street, Jeffreys Bay \,|\, +27 82 960 9353 \,|\, pauljeggelsdesigns.co.za}
\fancyfoot[R]{\small\thepage}
\renewcommand{\headrulewidth}{0.4pt}
\renewcommand{\footrulewidth}{0.4pt}

% Q&A formatting
\newcommand{\question}[1]{\subsection*{Q: #1}}
\newcommand{\answer}{\noindent\textbf{A:} }

\begin{document}

\begin{center}
{\LARGE\bfseries\color{pjdnavy} Paul Jeggels Designs}\\[6pt]
{\Large\color{pjdnavy} Surfboard Selection \& Customisation FAQ}\\[4pt]
{\small\color{gray} Hand-Shaped in Jeffreys Bay Since the 1980s}
\end{center}

\vspace{1em}

% ─────────────────────────────────────────────────────────────────────
\section{Choosing Your First Board}
% ─────────────────────────────────────────────────────────────────────

\question{I'm a complete beginner. What board should I start with?}
\answer If you're just learning, you want something with plenty of volume --- that means a board that's wide, thick, and long enough to keep you stable and help you catch waves easily. For most beginners, a longboard (around 8'0 to 9'0) or a generous funboard is the way to go. Paul will ask about your height, weight, and fitness level to get the sizing right. You don't need anything fancy --- you need something that paddles well and forgives your mistakes while you learn.

\question{What does ``volume'' mean, and why does it matter?}
\answer Volume is the total amount of space inside your board, measured in litres. More volume means more float and easier paddling. Less volume means a more responsive, performance-oriented board --- but it's harder to paddle and catch waves on. A beginner might ride a board with 50--70 litres. An experienced surfer on a shortboard might be down around 25--30 litres. Paul factors your weight, fitness, and surfing ability into the volume calculation so you get a board that works from your very first session.

\question{What's the difference between a shortboard, fish, longboard, and funboard?}
\answer Here's the simple version:
\begin{itemize}[leftmargin=1.5em]
    \item \textbf{Shortboard} (5'6--6'6): Built for speed, sharp turns, and powerful surfing. Best for experienced surfers who ride regularly.
    \item \textbf{Fish} (5'4--6'2): Wider and flatter than a shortboard. Quick, loose, and great in smaller waves. A favourite for intermediate to advanced surfers who want fun in everyday conditions.
    \item \textbf{Funboard / Mid-length} (6'6--8'0): The middle ground. Good paddle power, stable, and versatile. Works for intermediates and anyone who wants one board for everything.
    \item \textbf{Longboard} (8'0--10'0): Maximum stability and wave-catching ability. Smooth, flowing style. Perfect for beginners and anyone who loves cruisy surfing.
\end{itemize}
Paul shapes all of these. If you're not sure which is right, that's exactly what the consultation is for.

\question{I've been surfing for a couple of years. Should I go shorter?}
\answer Maybe --- but don't rush it. A lot of intermediate surfers drop volume too fast and end up frustrated. Paul's advice is to go down gradually. If you're comfortable on a 7'0, try a 6'6 next --- not a 5'10. The goal is to keep catching waves and having fun while your skills develop. Paul will look at where you're at and suggest the right step down for you.

% ─────────────────────────────────────────────────────────────────────
\section{Board Dimensions Explained}
% ─────────────────────────────────────────────────────────────────────

\question{What do all the numbers on a surfboard mean?}
\answer When you see something like 6'2 x 20 x 2.5, that's:
\begin{itemize}[leftmargin=1.5em]
    \item \textbf{Length} (6'2 = 6 feet 2 inches): Affects speed and how early you catch waves. Longer boards paddle faster.
    \item \textbf{Width} (20 inches): Affects stability. Wider = more stable but less responsive in turns.
    \item \textbf{Thickness} (2.5 inches): Affects float and paddle power. Thicker = more volume.
\end{itemize}
There's also \textbf{rocker} (the curve from nose to tail) and \textbf{concave} (the subtle channels on the bottom) which Paul dials in based on the waves you ride. You don't need to understand every detail --- Paul does. But if you're curious, he's happy to explain it all when you pick up your board.

\question{What is rocker and why does it matter?}
\answer Rocker is the curve of the board when you look at it from the side. More rocker in the nose stops you from nosediving on steep waves. More rocker in the tail makes the board looser and easier to turn. Flatter rocker means more speed but less manoeuvrability. For Jeffreys Bay's long, fast right-handers, Paul typically shapes a moderate rocker --- enough curve to handle the power, but flat enough to maintain speed down the line.

% ─────────────────────────────────────────────────────────────────────
\section{Custom vs Off-the-Rack}
% ─────────────────────────────────────────────────────────────────────

\question{Why should I get a custom board instead of buying one off the rack?}
\answer An off-the-rack board is designed for a generic surfer at a generic weight in generic waves. A custom board is designed for \emph{you} --- your body, your ability, and the actual waves you surf. Paul asks where you paddle out, how often you surf, what you struggle with, and what you want to improve. That conversation shapes the board. You can't get that from a factory. The difference is like wearing tailored clothes versus buying off the shelf --- both work, but one fits perfectly.

\question{Is a custom board worth it for a beginner?}
\answer It can be, yes. A beginner custom board isn't about performance features --- it's about getting the right volume and dimensions for your specific body so you progress faster. That said, if you're brand new, one of Paul's second-hand boards might be the smarter move. You'll get a properly shaped board at a lower price, and once you know what you like, you can order a custom shape that's dialled in for you.

\question{How long does a custom board last?}
\answer With proper care, a well-built custom board can last years. Paul uses quality materials --- Clark or Burford blanks, premium fibreglass, and UV-stable resin. The glass job (how thick the fibreglass is) affects durability. A heavier glass job lasts longer but adds weight. Paul will recommend the right balance based on how you surf and how you treat your gear.

% ─────────────────────────────────────────────────────────────────────
\section{Customisation Options}
% ─────────────────────────────────────────────────────────────────────

\question{What can I actually customise on my board?}
\answer Pretty much everything:
\begin{itemize}[leftmargin=1.5em]
    \item \textbf{Shape}: Length, width, thickness, nose and tail shape, rocker, concave, rails
    \item \textbf{Fin setup}: Single, twin, thruster (3 fins), quad (4 fins), or 5-fin (convertible)
    \item \textbf{Fin system}: FCS or Futures --- both work well, it's personal preference
    \item \textbf{Glass job}: Standard (4+4/4) for lighter weight, or heavier (6+4/6) for durability
    \item \textbf{Colour and finish}: Tints, sprays, airbrushed artwork, or clean white
    \item \textbf{Tail shape}: Squash, round, swallow, pin, diamond --- each changes how the board turns and holds in the wave
\end{itemize}
Paul will guide you through what matters for your surfing and what's purely cosmetic.

\question{Can I get custom artwork or a specific colour?}
\answer Absolutely. Paul does colour tints, resin swirls, and can arrange airbrushed designs. If you have something specific in mind --- a logo, a pattern, a particular colour --- bring it up during the consultation. Simple tints and sprays are included or minimal extra cost. Detailed airbrush work may add to the price depending on complexity.

% ─────────────────────────────────────────────────────────────────────
\section{Boards for Jeffreys Bay}
% ─────────────────────────────────────────────────────────────────────

\question{What board works best for Jeffreys Bay waves?}
\answer J-Bay is famous for its long, fast, hollow right-hand point breaks --- Supertubes being the crown jewel. For those waves, you want a board with good speed (moderate rocker, not too much flip in the nose), a responsive tail for linking turns on the open face, and enough volume to paddle into waves early. Paul has surfed every break in town for over 40 years, so when he shapes a board for J-Bay, he knows exactly what it needs. If you mainly surf Supers, Kitchen Windows, or Point, he'll tune the board for those specific waves.

\question{I don't surf J-Bay. Can Paul still shape for my local break?}
\answer Of course. Paul shapes boards for surfers all over South Africa and internationally. He'll ask you about your local conditions --- is it beach break or point break? Punchy and hollow, or slow and mellow? Sandy bottom or reef? That information tells him how to set the rocker, volume distribution, and fin setup. You don't have to surf J-Bay to get a board from Paul --- you just need to tell him where you surf.

\question{What board should I ride if I surf a mix of different spots?}
\answer If you surf beach breaks one day and points the next, Paul will shape something versatile --- typically a fish or a slightly fuller shortboard with moderate rocker. The idea is a board that handles a range of conditions without being too specialised for any single wave type. Think of it as your daily driver. Many of Paul's customers ride one all-rounder for 80\% of their sessions and have a second board for when conditions get really good.

\vfill
\begin{center}
\small\color{gray}
Paul Jeggels Designs --- 15 Dageraad Street, Jeffreys Bay, South Africa\\
Phone: +27 82 960 9353 --- Hand-shaped surfboards since the 1980s
\end{center}

\end{document}
